\documentclass{article}
\usepackage{graphicx}
\usepackage{amsfonts}
\usepackage{amsmath}
\usepackage{amsthm}
\usepackage{tikz}

\theoremstyle{definition}
\newtheorem{theorem}{Theorem}[section]
\newtheorem{corollary}[theorem]{Corollary}
\newtheorem{proposition}[theorem]{Proposition}
\newtheorem{lemma}[theorem]{Lemma}
\newtheorem{definition}[theorem]{Definition}

\theoremstyle{remark}
\newtheorem{remark}[theorem]{Remark}
\newtheorem{example}[theorem]{Example}
\newtheorem{claim}[theorem]{Claim}

\title{Mean Value Theorem}
\author{Sarah Liu}
\date{January 2026}

\begin{document}

\maketitle

\begin{theorem}[Fermat's Theorem]
    Let $f$ be defined on $(a,b)$ and let $c\in(a,b)$ be such that $f(c)\geq f(x)$ for all $x\in(a,b)$. Then either $f^{'}(c)=0$ or $f^{'}(c)$ does not exist.
\end{theorem}

\begin{theorem}[Rolle's Theorem]
    Let $f$ be continuous on $[a,b]$ and differentiable on $(a,b)$, with $f(a)=f(b)$. Then there exists $c\in(a,b)$ such that $f^{'}(c)=0$.
\end{theorem}

\begin{proof}
    If $f$ is constant on $[a,b]$, then its derivative is zero and so any $c\in(a,b)$ satisfies the conclusion of the theorem.
    So we assume that $f$ is not constant on $[a,b]$.
    By the Extreme Value Theorem, $f$ attains an absolute maximum and an absolute minimum on $[a,b]$.
    Since $f$ is not constant, at least one of these absolute extrema must occur at $c\in(a,b)$.
    Then by Fermat's Theorem, we know that $f^{'}(c)=0$.
\end{proof}

\begin{theorem}[Cauchy's Mean Value Theorem]
    Let $f$ and $g$ be continuous on $[a,b]$ and differentiable on $(a,b)$. Then there exists $c\in(a,b)$ such that 
    \[
    f^{'}(c)[g(b)-g(a)]=g^{'}(c)[f(b)-f(a)]
    \]
\end{theorem}

\begin{proof}
    If $g(a)=g(b)$, then by Rolle's Theorem, the result holds.
    We assume that $g(a)\neq g(b)$.
    We consider the function
    \[
    h(x)=f(x)-\frac{f(b)-f(a)}{g(b)-g(a)}g(x).
    \]
    Then $h$ is continuous on $[a,b]$ and differentiable on $(a,b)$.
    And $h(a)=h(b)$.
    By Rolle's Theorem, there is $c\in(a,b)$ such that $h^{'}(c)=0$.
    Finally, $h^{'}(c)=0$ gives $f^{'}(c)[g(b)-g(a)]=g^{'}(c)[f(b)-f(a)]$.
\end{proof}

\end{document}